\section{Limitations}
\label{sec:limitations}

This study covers a single dataset, language, and pair of source
platforms: EDOS's English-language Reddit and Gab posts
\cite{kirk2023semeval}. Sexist language is expressed differently across
languages and platforms (Sec.~\ref{sec:related:datasets} discusses SWSR
\cite{jiang2021swsr}, the closest non-English counterpart), so nothing
here establishes that the same prompt-design effects, the same baseline
gap, or the same safety-refusal pattern would reproduce on a different
dataset. The study also targets only EDOS's binary sexist/not-sexist
subtask directly; the dataset's finer four-category taxonomy is used
solely for the post-hoc subtype error analysis
(Sec.~\ref{sec:results:subtype}), never as a training or prediction
target in its own right.

The prompt-design grid is a curated 18 of 32 possible component
combinations (Sec.~\ref{sec:methodology:prompts}), not an exhaustive
search, and every component's wording is hand-designed rather than
automatically optimized or drawn from a cited theoretical taxonomy; this
includes the \textbf{aspects} component's seven linguistic markers, which
reflect our own judgment of what sexist language looks like rather
than an established sociolinguistic framework. A negative result
within this space, most
prompted variants scoring below the classical baseline, is evidence
about the designs actually tested, not proof that no possible prompt
formulation could close the gap. Similarly, the few-shot demonstrations
for every main-grid variant are drawn once, with a fixed seed, from the
training split. Sec.~\ref{sec:methodology:sweeps}'s order, count, and
demonstration-seed sweeps test how sensitive each model's own best
variant is to arrangement, count, and which specific examples are drawn,
respectively (Sec.~\ref{sec:results:sensitivity},
Sec.~\ref{sec:results:seed-sensitivity}), but only for that one variant
per model, not the full grid; the seed sweep itself covers three
alternate draws per model, enough to surface phi3's large,
model-specific sensitivity to selection, but not an exhaustive
characterization of the selection space.

The model roster is six open-weight, instruction-tuned models in a
3.8--9B parameter range (Table~\ref{tab:models}). No larger open model
and no closed, commercial API model, the kind increasingly used in
production moderation systems, is evaluated, so neither the
prompt-design findings nor the baseline comparison necessarily extend to
that different and increasingly common deployment setting. Every
inference-cost figure in this study (Table~\ref{tab:cost}) is specific
to this study's own hardware and batching choices;
the direction of the cost argument in
Sec.~\ref{sec:discussion:qwen3-cost}, that qwen3's win over the
classical baseline is expensive, is likely to hold generally, but its
exact multiplier is not.

This study compares against a classical machine-learning baseline and
two fine-tuned neural baselines, but not against human inter-annotator
agreement on the same test items. Without that reference point, it is
not possible to say how close, or how far, any of this study's numbers
sit from whatever ceiling human disagreement on sexism labeling itself
imposes.

Sec.~\ref{sec:results:llama-refusal} and Table~\ref{tab:llama-cot} show
that llama3.1's elevated \texttt{fail\_ratio} on chain-of-thought
variants splits into two distinct causes depending on whether aspects is
active: explicit safety refusal when it is not, and token-budget
truncation, the model's own step-by-step reasoning running past this
study's fixed 220-token chain-of-thought budget, when it is. That budget
was calibrated for ``brief prose reasoning'' across the five
non-reasoning models generally
(Sec.~\ref{sec:experimental:generation}); it is demonstrably too short
for at least this one model under at least this one prompt
configuration. Whether the same refusal mechanism, at any scale,
generalizes to the other five models has since been checked
systematically rather than spot-checked: every chain-of-thought cell
with at least one failure across mistral, phi3, qwen2.5, gemma2, and
qwen3 (47 cells) was classified by the same refusal-pattern rule used
for Table~\ref{tab:llama-cot}, reading explicit-refusal
\emph{prevalence} as a share of the full 4{,}000-example test split
rather than a share of each cell's own failures, since the two can
diverge sharply when \texttt{fail\_ratio} itself is small. On that
measure, explicit refusal never exceeds 1.05\% of the test split for any
of the five (phi3's V13, the highest rate observed outside llama3.1;
mistral peaks at 0.30\%, gemma2 at 0.23\%, qwen3 at 0.15\%, and qwen2.5
shows none). Every one of llama3.1's own ten other chain-of-thought
cells sits well above that entire range, from 4.82\% at its lowest to
48.1\% at V8. Llama3.1's safety-refusal interference
(Sec.~\ref{sec:results:llama-refusal}) is, on this evidence, a property
of that model specifically, not a shared chain-of-thought failure mode
this study's prompt design triggers generally.

Sec.~\ref{sec:results:auc}'s AUC comparison, and
Sec.~\ref{sec:results:abstention}'s confidence-based abstention policy
built on the same score, both cover only the seven non-chain-of-thought
variants. The eleven chain-of-thought variants, including qwen3's own
best-F1 variant (V8), have no single decision token this method can
score, so neither AUC nor an abstention gain exists for them at all. The
abstention gain itself is also not free or uniform: Sec.~\ref{sec:results:abstention}
shows one model (phi3) buying a large F1 improvement only by abstaining
on over a third of the test split, while four of the other five models
gain comparatively little at a much smaller coverage cost, so the same
band width does not represent the same real-world tradeoff across
models.

Generation throughout this study is greedy and seeded, which removes
sampling randomness but not every source of nondeterminism: kernel-level
differences across GPU architectures, driver versions, and library
versions can still produce different outputs from bit-identical inputs
on different hardware. Each of five models' best variant was rerun on two cluster
nodes with the same GPU architecture (GA100, compute capability 8.0) but
different memory configurations, a 40GB-per-card node and an
80GB-per-card node: every one of 4{,}000 test examples produced
byte-identical raw output and an identical parsed decision on both
nodes, for all five models, confirming this study's pipeline is
bit-for-bit deterministic across nodes and memory configurations within
one GPU architecture. The two-run reproduction of llama3.1's V14 refusal
rate (Sec.~\ref{sec:results:llama-refusal}) is a special case of the
same result. qwen3, whose batch-128 footprint does not fit the 40GB
node's cards, was instead rerun on the 80GB node's A100 cards, again
matching Table~\ref{tab:main-grid}'s own V8 figure exactly. What this
study can report is strong, directly tested evidence for bit-identical
determinism across nodes and memory configurations within one GPU
architecture; whether the same holds across different GPU architectures
remains untested.
